% ═══════════════════════════════════════════════════════════════════════════
% QR-LP-03: Elektrische Leistung & Energie
% Physik Klasse 8 | Greenhouse Schools Rostock
% ═══════════════════════════════════════════════════════════════════════════

\documentclass[a4paper,11pt]{article}

\usepackage[utf8]{inputenc}
\usepackage[T1]{fontenc}
\usepackage{helvet}
\renewcommand{\familydefault}{\sfdefault}

\usepackage[margin=20mm]{geometry}
\usepackage{tikz}
\usepackage{xcolor}
\usepackage{qrcode}
\usepackage{hyperref}

% ═══════════════════════════════════════════════════════════════════════════
% FARBEN
% ═══════════════════════════════════════════════════════════════════════════

\definecolor{physik}{HTML}{2c5aa0}
\definecolor{chemie}{HTML}{2a7a4b}
\definecolor{informatik}{HTML}{b35c00}

% ═══════════════════════════════════════════════════════════════════════════
% KONFIGURATION
% ═══════════════════════════════════════════════════════════════════════════

% Fach (physik | chemie | informatik)
\newcommand{\fachid}{physik}
\colorlet{fachfarbe}{\fachid}

% Titel des Lernpfads
\newcommand{\lptitel}{Elektrische Leistung \& Energie}

% URL zum Online-Lernpfad
\newcommand{\lpurl}{https://devbox42.github.io/sciencesim/physik/kl08/03-leistung/LP-03-elektrische-leistung.html}

% Kurz-URL zur Anzeige
\newcommand{\lpurlkurz}{devbox42.github.io/.../LP-03-elektrische-leistung.html}

% ═══════════════════════════════════════════════════════════════════════════
% DOKUMENT
% ═══════════════════════════════════════════════════════════════════════════

\pagestyle{empty}

\begin{document}

\begin{center}

% Titel
\vspace*{10mm}
{\fontsize{28}{34}\selectfont\bfseries\color{fachfarbe}Lernpfad: \lptitel}

\vspace{15mm}

% QR-Code mit Rahmen (klickbar)
\href{\lpurl}{\fcolorbox{fachfarbe}{white}{\qrcode[height=100mm]{\lpurl}}}

\vspace{12mm}

% URL
{\large\color{gray}\lpurlkurz}

\vspace{8mm}

% Hinweis
{\Large\color{gray}Scanne den QR-Code mit deinem Gerät}

\vspace{10mm}

% Info-Box
\fcolorbox{fachfarbe}{fachfarbe!10}{%
    \parbox{0.7\textwidth}{%
        \centering
        \vspace{3mm}
        {\large Arbeite in deinem Tempo durch den Lernpfad.}\\[2mm]
        Nutze das Arbeitsblatt zum Notieren.
        \vspace{3mm}
    }%
}

\end{center}

\end{document}
